% !TEX program = xelatex
\documentclass[11pt,a4paper]{article}

\usepackage[UTF8,fontset=none]{ctex}
\setCJKmainfont[BoldFont=Heiti SC,ItalicFont=Songti SC]{Songti SC}
\setCJKsansfont{Heiti SC}
\setCJKmonofont{Heiti SC}

\usepackage{amsmath,amssymb,amsfonts,mathtools}
\usepackage{booktabs,array,tabularx}
\usepackage[margin=1.9cm]{geometry}
\usepackage[dvipsnames]{xcolor}
\usepackage[most]{tcolorbox}
\usepackage[hidelinks]{hyperref}
\usepackage{enumitem}
\usepackage{fancyhdr}

\setlength{\parindent}{0em}
\setlength{\parskip}{0.48em}
\setlist[itemize]{leftmargin=*,topsep=2pt,itemsep=2pt}
\setlist[enumerate]{leftmargin=*,topsep=2pt,itemsep=2pt}
\allowdisplaybreaks

\pagestyle{fancy}
\fancyhf{}
\lhead{\small Statistical Modelling and Inference}
\rhead{\small 第 1--2 部分详细讲解}
\cfoot{\thepage}

\hypersetup{
  pdftitle={SMI Parts 1 and 2 Detailed Explanation},
  pdfauthor={Prepared from SMI Final Review 2026}
}

\newcommand{\E}{\operatorname{E}}
\newcommand{\Var}{\operatorname{Var}}
\newcommand{\Cov}{\operatorname{Cov}}
\newcommand{\Bias}{\operatorname{Bias}}
\newcommand{\MSE}{\operatorname{MSE}}
\newcommand{\SE}{\operatorname{SE}}
\newcommand{\iid}{\overset{\mathrm{iid}}{\sim}}

\newtcolorbox{goalbox}[1][]{
  colback=gray!5!white,colframe=gray!65!black,
  fonttitle=\bfseries,title=#1,breakable
}
\newtcolorbox{conceptbox}[1][]{
  colback=blue!4!white,colframe=blue!60!black,
  fonttitle=\bfseries,title=#1,breakable
}
\newtcolorbox{intuitionbox}[1][]{
  colback=cyan!4!white,colframe=cyan!60!black,
  fonttitle=\bfseries,title=#1,breakable
}
\newtcolorbox{methodbox}[1][]{
  colback=green!4!white,colframe=green!55!black,
  fonttitle=\bfseries,title=#1,breakable
}
\newtcolorbox{examplebox}[1][]{
  colback=orange!6!white,colframe=orange!75!black,
  fonttitle=\bfseries,title=#1,breakable
}
\newtcolorbox{proofbox}[1][]{
  colback=violet!4!white,colframe=violet!60!black,
  fonttitle=\bfseries,title=#1,breakable
}
\newtcolorbox{warnbox}[1][]{
  colback=red!4!white,colframe=red!60!black,
  fonttitle=\bfseries,title=#1,breakable
}
\newtcolorbox{summarybox}[1][]{
  colback=yellow!7!white,colframe=Goldenrod!75!black,
  fonttitle=\bfseries,title=#1,breakable
}

\title{\textbf{Statistical Modelling and Inference}\\[0.3em]
第 1、2 部分详细易懂讲解}
\author{对应《SMI 2026 期末复习讲义》第 1--2 章}
\date{June 2026}

\begin{document}
\maketitle

\begin{goalbox}[阅读目标]
完成本讲解后，应能够：
\begin{enumerate}
  \item 区分总体参数、统计量、估计量与估计值；
  \item 解释抽样分布和标准误，而不是只背公式；
  \item 面对均值或比例推断题时，先判断参数与研究设计，再选择方法；
  \item 判断估计量是否无偏，计算 Bias、Variance 与 MSE；
  \item 证明 MSE 的偏差--方差分解；
  \item 构造无偏估计量，并合并多个无偏估计量；
  \item 解释样本均值为何无偏，以及样本方差为何使用分母 $n-1$。
\end{enumerate}
\end{goalbox}

\tableofcontents
\newpage

% ============================================================
\section{第 1 部分：全课程主线与做题决策}
% ============================================================

\subsection{统计推断究竟在做什么}

统计问题通常具有一个根本矛盾：
\[
  \boxed{\text{我们关心整个总体，但只能观察总体中的一个样本。}}
\]

例如，我们可能想知道全校学生的平均睡眠时长 $\mu$。真实的 $\mu$ 在调查前就已经由总体决定，但我们无法询问所有学生，只能随机抽取 $n$ 人并计算样本均值 $\bar X$。统计推断所做的事情，就是利用样本中可观察的信息，对不可观察的总体参数作出有不确定性说明的结论。

\[
  \boxed{
  \text{总体参数未知}
  \quad \xleftarrow{\quad\text{抽样分布}\quad}\quad
  \text{样本统计量可观察}.
  }
\]

\begin{conceptbox}[参数固定，样本与统计量随机]
在频率学派的框架中：
\begin{itemize}
  \item 总体参数 $\theta$ 是\textbf{固定但未知}的数；
  \item 随机样本 $X_1,\ldots,X_n$ 在抽样前是随机变量；
  \item 由样本计算的统计量 $T=g(X_1,\ldots,X_n)$ 也在抽样前是随机变量；
  \item 数据观察完成后，统计量取得一个具体数值。
\end{itemize}
因此，“$\theta$ 的概率是多少”通常不是频率学派问题；我们研究的是估计方法在重复抽样中的表现。
\end{conceptbox}

\subsection{总体、样本、参数与统计量}

\begin{center}
\renewcommand{\arraystretch}{1.25}
\begin{tabularx}{\textwidth}{p{3.1cm}p{3.5cm}X}
\toprule
概念 & 含义 & 例子\\
\midrule
总体 population & 研究所关心的全部对象或概率机制
& 全校学生；某机器未来生产的所有零件\\
样本 sample & 从总体中实际观察的一组数据
& 随机抽取的 50 名学生的睡眠时长\\
参数 parameter & 描述总体的固定未知数
& 总体均值 $\mu$、方差 $\sigma^2$、比例 $p$\\
统计量 statistic & 只由样本计算、不含未知参数的函数
& $\bar X$、$S^2$、$\hat p$\\
\bottomrule
\end{tabularx}
\end{center}

\begin{warnbox}[统计量不能含未知参数]
例如
\[
  \bar X=\frac1n\sum_{i=1}^nX_i
\]
是统计量，因为它只依赖样本；$\bar X-\mu$ 含未知参数，因此不是统计量。
\end{warnbox}

\subsection{Estimator、Estimate 与 Estimand}

这三个词很相似，但扮演不同角色：

\begin{itemize}
  \item \textbf{Estimand（被估计对象）}：真正想知道的总体量，例如 $\mu$ 或 $p$；
  \item \textbf{Estimator（估计量）}：用于估计参数的随机规则，例如 $\bar X$；
  \item \textbf{Estimate（估计值）}：把观测数据代入估计量后得到的具体数值，例如 $\bar x=7.2$。
\end{itemize}

\begin{examplebox}[同一估计量，不同样本会产生不同估计值]
设总体均值为未知参数 $\mu$，使用估计量
\[
  \hat\mu=\bar X.
\]
第一次抽到数据 $(5,7,9)$，得到
\[
  \bar x=7.
\]
第二次抽到数据 $(4,6,8)$，得到
\[
  \bar x=6.
\]
两次使用的是同一个\textbf{估计量} $\bar X$，但产生了不同的\textbf{估计值}。这正是估计量具有抽样分布的原因。
\end{examplebox}

\subsection{抽样分布：连接样本与总体的桥梁}

想象我们从同一个总体中反复抽取大小均为 $n$ 的随机样本。每次都计算同一个统计量，例如 $\bar X$。这些可能的 $\bar X$ 值形成的概率分布，称为 $\bar X$ 的\textbf{抽样分布（sampling distribution）}。

\begin{intuitionbox}[不要把三个分布混在一起]
\begin{itemize}
  \item \textbf{总体分布}：单个个体 $X$ 如何变化；
  \item \textbf{样本的经验分布}：这一次实际数据如何分布；
  \item \textbf{统计量的抽样分布}：重复抽样时，$\bar X$、$\hat p$ 等如何变化。
\end{itemize}
置信区间和假设检验主要依赖第三种分布。
\end{intuitionbox}

若 $X_1,\ldots,X_n$ 独立同分布，且
\[
  \E(X_i)=\mu,
  \qquad
  \Var(X_i)=\sigma^2,
\]
则
\[
  \E(\bar X)=\mu,
  \qquad
  \Var(\bar X)=\frac{\sigma^2}{n}.
\]
因此 $\bar X$ 的标准差为
\[
  \boxed{\SE(\bar X)=\frac{\sigma}{\sqrt n}.}
\]

\begin{conceptbox}[标准差与标准误的区别]
\begin{itemize}
  \item $\sigma$ 或 $s$ 描述\textbf{个体观测}之间的差异；
  \item $\SE(\bar X)$ 描述\textbf{样本均值}在重复抽样中的波动。
\end{itemize}
样本量变大不会让个体差异自动消失，但会让样本均值更稳定。由于
\[
  \SE(\bar X)=\frac{\sigma}{\sqrt n},
\]
若希望标准误缩小一半，样本量需要扩大到原来的四倍。
\end{conceptbox}

\subsection{为什么推断题必须先定义参数}

公式选择取决于你究竟想回答什么问题。以下问题看起来都在比较数据，但目标参数不同：

\begin{center}
\renewcommand{\arraystretch}{1.25}
\begin{tabularx}{\textwidth}{p{6.0cm}p{3.0cm}X}
\toprule
研究问题 & 目标参数 & 典型估计量\\
\midrule
一个总体的平均值是多少？ & $\mu$ & $\bar X$\\
两个独立总体的平均值相差多少？ & $\mu_1-\mu_2$ & $\bar X_1-\bar X_2$\\
同一批对象干预前后平均变化多少？ & $\mu_d$ & $\bar D$\\
总体中成功或支持的比例是多少？ & $p$ & $\hat p=X/n$\\
两个总体比例相差多少？ & $p_1-p_2$ & $\hat p_1-\hat p_2$\\
\bottomrule
\end{tabularx}
\end{center}

\begin{methodbox}[定义参数的标准写法]
不要只写“检验两个均值是否相同”。更完整的写法是：
\[
  \mu_1=\text{总体 1 的平均响应},
  \qquad
  \mu_2=\text{总体 2 的平均响应},
\]
目标参数为 $\mu_1-\mu_2$。参数定义清楚后，假设、标准误、置信区间和语境结论才不会混乱。
\end{methodbox}

\subsection{研究设计优先于公式：独立样本还是配对样本}

两组数据并不自动意味着两个独立样本。关键问题是：\textbf{两组观测之间是否存在一一对应关系？}

\begin{itemize}
  \item 不同学生被随机分入两种教学方法：通常是两个独立样本；
  \item 同一批学生课程前后各测一次：配对样本；
  \item 双胞胎按对分到两个处理组：配对样本；
  \item 从两个互不重叠的人群分别抽样：通常是独立样本。
\end{itemize}

配对数据应先定义差值
\[
  D_i=X_{1i}-X_{2i},
\]
再把问题转化为对单总体均值 $\mu_d$ 的推断。若错误地把配对数据当作独立样本，会忽略配对关系携带的信息，并使用错误的标准误。

\begin{examplebox}[设计判断]
某研究记录 30 名患者服药前后的血压。虽然有“服药前”和“服药后”两列数据，但每个服药后观测都对应同一患者的服药前观测，因此这是配对设计。目标参数是
\[
  \mu_d=\E(\text{服药后血压}-\text{服药前血压}).
\]
正确做法是分析 30 个差值，而不是把两列当作两个独立样本。
\end{examplebox}

\subsection{\texorpdfstring{$Z$ 与 $t$：为什么总体方差未知时会改变参考分布}{Z 与 t：为什么总体方差未知时会改变参考分布}}

若总体正态且总体标准差 $\sigma$ 已知，则
\[
  Z
  =
  \frac{\bar X-\mu}{\sigma/\sqrt n}
  \sim N(0,1).
\]

实际中 $\sigma$ 常常未知，只能用样本标准差 $S$ 替代：
\[
  T
  =
  \frac{\bar X-\mu}{S/\sqrt n}.
\]
分母 $S$ 本身也是随机的，因此统计量比标准正态具有更厚的尾部。在总体正态条件下，
\[
  \boxed{
  T\sim t_{n-1}.
  }
\]

\begin{intuitionbox}[为什么样本量大时 $t$ 越来越接近 $Z$]
样本量较小时，$S$ 对 $\sigma$ 的估计不够稳定，替换 $\sigma$ 带来较多额外不确定性，因此 $t$ 分布尾部更厚。样本量增大后，$S$ 越来越稳定，额外不确定性减小，$t$ 分布逐渐接近标准正态。
\end{intuitionbox}

\subsection{方法选择的逐步决策}

\begin{methodbox}[任何均值或比例推断题的决策顺序]
\begin{enumerate}
  \item \textbf{定义目标参数}：$\mu$、$\mu_1-\mu_2$、$\mu_d$、$p$ 或 $p_1-p_2$？
  \item \textbf{识别设计}：单样本、两个独立样本，还是配对样本？
  \item \textbf{识别信息与条件}：方差是否已知？样本量大小？是否可假设正态或使用大样本近似？
  \item \textbf{选择标准误}：标准误必须与目标估计量相匹配。
  \item \textbf{选择参考分布}：标准正态 $Z$ 还是某个自由度的 $t$？
  \item \textbf{计算并解释}：最后才代数值，并用总体参数的语言写结论。
\end{enumerate}
\end{methodbox}

\begin{center}
\renewcommand{\arraystretch}{1.23}
\begin{tabularx}{\textwidth}{p{2.3cm}p{4.0cm}X}
\toprule
目标 & 主要情形 & 核心方法\\
\midrule
$\mu$ & $\sigma$ 已知 & $Z$，标准误 $\sigma/\sqrt n$\\
$\mu$ & $\sigma$ 未知、总体正态 & $t_{n-1}$，标准误 $s/\sqrt n$\\
$\mu_1-\mu_2$ & 独立、小样本、正态、等方差 & pooled $t$\\
$\mu_1-\mu_2$ & 独立、不假设等方差 & Welch $t$；按课程或题目要求使用\\
$\mu_d$ & 配对数据 & 对差值做单样本 $t$\\
$p$ & 单总体比例、大样本 & 正态近似；CI 与检验标准误不同\\
$p_1-p_2$ & 两个独立比例、大样本 & CI 不合并；等比例检验时合并\\
\bottomrule
\end{tabularx}
\end{center}

\subsection{条件检查不是装饰}

推断公式来自特定概率模型。若条件明显不满足，公式给出的置信水平或错误概率就可能失真。

\begin{itemize}
  \item \textbf{随机性或代表性}：样本应能代表目标总体；
  \item \textbf{独立性}：一个观测不应决定另一个观测，除非明确建模配对结构；
  \item \textbf{正态性}：小样本均值 $t$ 推断常依赖总体近似正态；
  \item \textbf{大样本条件}：比例正态近似需要有足够的成功与失败数；
  \item \textbf{等方差条件}：pooled $t$ 需要两个总体方差相等。
\end{itemize}

\begin{warnbox}[“样本量大”不能修复所有问题]
大样本有助于正态近似，但不能自动修复严重选择偏差、错误配对、观测依赖或测量错误。统计方法只能量化模型中承认的不确定性，不能消除糟糕研究设计带来的系统问题。
\end{warnbox}

\subsection{三个方法选择示例}

\begin{examplebox}[示例 1：单总体均值]
随机抽取 16 个来自正态总体的观测，总体方差未知，希望估计总体均值 $\mu$。

\textbf{判断：}
\[
  \text{单样本}
  \quad+\quad
  \sigma\text{ 未知}
  \quad+\quad
  \text{总体正态}
  \quad\Longrightarrow\quad
  t_{15}.
\]
标准误为 $s/\sqrt{16}$。
\end{examplebox}

\begin{examplebox}[示例 2：两个均值]
分别从两个独立正态总体抽取小样本，题目明确说明总体方差相同，希望研究 $\mu_1-\mu_2$。

\textbf{判断：}
\[
  \text{两个独立样本}
  \quad+\quad
  \text{小样本正态}
  \quad+\quad
  \text{等方差}
  \quad\Longrightarrow\quad
  \text{pooled }t.
\]
应先计算合并方差 $s_p^2$，自由度为 $n_1+n_2-2$。
\end{examplebox}

\begin{examplebox}[示例 3：总体比例]
随机调查 200 人，其中 116 人支持某提议，希望估计总体支持比例 $p$。

点估计为
\[
  \hat p=\frac{116}{200}=0.58.
\]
比例置信区间使用
\[
  \sqrt{\frac{\hat p(1-\hat p)}{n}}.
\]
若检验 $H_0:p=0.5$，则检验标准误应在原假设下计算：
\[
  \sqrt{\frac{0.5(1-0.5)}{n}}.
\]
两者不能混用。
\end{examplebox}

\begin{summarybox}[第 1 部分的核心]
统计推断不是“看到题型就找公式”，而是：
\[
  \boxed{
  \text{定义参数}
  \longrightarrow
  \text{识别设计}
  \longrightarrow
  \text{确定抽样分布与标准误}
  \longrightarrow
  \text{计算并解释}.
  }
\]
只要前面三步正确，公式通常自然出现；若前面判断错误，即使计算完全正确，结论仍可能无效。
\end{summarybox}

% ============================================================
\section{第 2 部分：无偏估计、偏差与均方误差}
% ============================================================

\subsection{评价估计量：中心位置与波动程度}

估计量 $\hat\theta$ 是随机变量。评价它时，至少要问两个问题：
\begin{enumerate}
  \item 它在重复抽样中的平均位置是否接近真实参数 $\theta$？
  \item 它在不同样本之间是否稳定？
\end{enumerate}

第一个问题对应\textbf{偏差（bias）}，第二个问题对应\textbf{方差（variance）}。

\subsection{偏差与无偏性}

定义
\[
  \boxed{
  \Bias(\hat\theta)
  =
  \E(\hat\theta)-\theta.
  }
\]

\begin{itemize}
  \item $\Bias(\hat\theta)>0$：估计量长期平均偏高；
  \item $\Bias(\hat\theta)<0$：估计量长期平均偏低；
  \item $\Bias(\hat\theta)=0$：估计量无偏。
\end{itemize}

\[
  \boxed{
  \hat\theta\text{ 无偏}
  \quad\Longleftrightarrow\quad
  \E(\hat\theta)=\theta.
  }
\]

\begin{intuitionbox}[无偏不代表每次都估得准确]
若 $\hat\theta$ 有时严重高估、有时严重低估，只要长期平均恰好等于 $\theta$，它仍然无偏。因此无偏性只描述估计量分布的\textbf{中心位置}，不描述分布有多分散。
\end{intuitionbox}

\begin{examplebox}[一个无偏但可能很不稳定的估计量]
设 $X_1,\ldots,X_n$ 的均值均为 $\mu$。则
\[
  T=X_1
\]
也是 $\mu$ 的无偏估计量，因为
\[
  \E(T)=\E(X_1)=\mu.
\]
但它只使用了第一个观测，方差为 $\sigma^2$。相比之下，
\[
  \bar X=\frac1n\sum X_i
\]
同样无偏，但方差仅为 $\sigma^2/n$。因此“无偏”不足以判断估计量是否优秀。
\end{examplebox}

\subsection{均方误差 MSE：总体准确性的度量}

估计误差为
\[
  \hat\theta-\theta.
\]
若直接取期望，正负误差可能互相抵消。因此使用平方误差：
\[
  \boxed{
  \MSE(\hat\theta)
  =
  \E\bigl[(\hat\theta-\theta)^2\bigr].
  }
\]

MSE 越小，表示估计量通常越接近真实参数。平方还会对较大的错误施加更重惩罚。

\subsection{MSE 的偏差--方差分解}

最重要的恒等式是
\[
  \boxed{
  \MSE(\hat\theta)
  =
  \Var(\hat\theta)
  +
  \bigl[\Bias(\hat\theta)\bigr]^2.
  }
\]

\begin{proofbox}[逐步证明]
首先在估计误差中加减 $\E(\hat\theta)$：
\[
\begin{aligned}
  \hat\theta-\theta
  &=
  \hat\theta-\E(\hat\theta)
  +
  \E(\hat\theta)-\theta\\
  &=
  \underbrace{\hat\theta-\E(\hat\theta)}_{\text{随机波动}}
  +
  \underbrace{\Bias(\hat\theta)}_{\text{系统偏差}}.
\end{aligned}
\]
记 $B=\Bias(\hat\theta)$。平方：
\[
  (\hat\theta-\theta)^2
  =
  \bigl(\hat\theta-\E\hat\theta\bigr)^2
  {}+2B\bigl(\hat\theta-\E\hat\theta\bigr)
  {}+B^2.
\]
两边取期望：
\[
\begin{aligned}
  \MSE(\hat\theta)
  &=
  \E\left[\bigl(\hat\theta-\E\hat\theta\bigr)^2\right]
  {}+2B\,\E\left[\hat\theta-\E\hat\theta\right]
  {}+B^2\\
  &=
  \Var(\hat\theta)+2B(0)+B^2.
\end{aligned}
\]
因此
\[
  \MSE(\hat\theta)=\Var(\hat\theta)+B^2.
\]
交叉项消失的根本原因是：随机波动
$\hat\theta-\E(\hat\theta)$ 的期望恒为零。
\end{proofbox}

\begin{conceptbox}[分解的含义]
\[
  \text{总体误差}
  =
  \text{随机不稳定性}
  +
  \text{系统偏离程度的平方}.
\]
若估计量无偏，则
\[
  \MSE(\hat\theta)=\Var(\hat\theta).
\]
若估计量有偏，比较优劣时不能只比较方差，也不能只比较偏差；必须同时考虑两者。
\end{conceptbox}

\subsection{有偏估计量可能具有更小的 MSE}

\begin{examplebox}[偏差--方差权衡]
设 $\bar X$ 是 $\mu$ 的无偏估计量，并且
\[
  \Var(\bar X)=4.
\]
考虑缩小后的估计量
\[
  T=0.8\bar X.
\]
当真实参数 $\mu=2$ 时，
\[
  \E(T)=0.8\mu=1.6,
  \qquad
  \Bias(T)=-0.4.
\]
同时
\[
  \Var(T)=0.8^2\Var(\bar X)=0.64(4)=2.56.
\]
因此
\[
  \MSE(T)=2.56+(-0.4)^2=2.72.
\]
而
\[
  \MSE(\bar X)=\Var(\bar X)=4.
\]
虽然 $T$ 有偏，但它通过显著降低方差取得了更小的 MSE。这个例子说明：\textbf{无偏是重要性质，但不是唯一标准。}
\end{examplebox}

\begin{warnbox}[MSE 可能依赖未知参数]
上例中
\[
  \MSE(0.8\bar X)
  =
  0.64\Var(\bar X)+0.04\mu^2,
\]
所以有偏估计量是否优于无偏估计量，可能取决于真实参数 $\mu$。比较时必须说明在哪些参数条件下成立。
\end{warnbox}

\subsection{常用期望、方差与协方差规则}

以下规则是无偏估计题的主要计算工具：
\[
  \E(aX+b)=a\E(X)+b,
\]
\[
  \Var(aX+b)=a^2\Var(X),
\]
\[
  \Var(aX+bY)
  =
  a^2\Var(X)
  +
  b^2\Var(Y)
  +
  2ab\Cov(X,Y).
\]

若 $X,Y$ 独立，则
\[
  \Cov(X,Y)=0.
\]
但反方向一般不成立：协方差为零只表示无线性关联，未必表示完全独立。

\subsection{通过线性修正构造无偏估计量}

若某估计量满足
\[
  \E(T)=a\theta+b,
  \qquad a\ne0,
\]
则可以逆向消除缩放和偏移：
\[
  \boxed{
  T^*=\frac{T-b}{a}.
  }
\]
验证：
\[
  \E(T^*)
  =
  \frac{\E(T)-b}{a}
  =
  \frac{a\theta+b-b}{a}
  =
  \theta.
\]

\begin{methodbox}[构造无偏估计量的答题步骤]
\begin{enumerate}
  \item 先计算原估计量 $T$ 的期望；
  \item 将期望整理为 $a\theta+b$；
  \item 对 $T$ 做相反的平移和缩放，得到 $(T-b)/a$；
  \item 最后再取一次期望，验证结果确实等于 $\theta$。
\end{enumerate}
\end{methodbox}

\begin{examplebox}[线性修正]
若
\[
  \E(T)=2\theta+3,
\]
则
\[
  T^*=\frac{T-3}{2}
\]
是 $\theta$ 的无偏估计量，因为
\[
  \E(T^*)=\frac{2\theta+3-3}{2}=\theta.
\]
\end{examplebox}

\begin{examplebox}[缩放后的样本均值]
设
\[
  T=\frac{n-1}{n}\bar X,
  \qquad
  \E(\bar X)=\mu.
\]
则
\[
  \E(T)=\frac{n-1}{n}\mu,
\]
所以
\[
  \Bias(T)
  =
  -\frac{\mu}{n}.
\]
将其乘以 $n/(n-1)$ 即可消除偏差：
\[
  T^*
  =
  \frac{n}{n-1}T
  =
  \bar X.
\]
\end{examplebox}

\subsection{合并两个无偏估计量}

设 $\hat\theta_1,\hat\theta_2$ 都是 $\theta$ 的无偏估计量。考虑
\[
  \hat\theta_a
  =
  a\hat\theta_1+(1-a)\hat\theta_2.
\]
因为权重和为 1，
\[
\begin{aligned}
  \E(\hat\theta_a)
  &=
  a\E(\hat\theta_1)+(1-a)\E(\hat\theta_2)\\
  &=
  a\theta+(1-a)\theta\\
  &=
  \theta.
\end{aligned}
\]
因此任意固定 $a$ 都保持无偏。接下来应选择使方差最小的 $a$。

若
\[
  \Var(\hat\theta_1)=v_1,
  \quad
  \Var(\hat\theta_2)=v_2,
  \quad
  \Cov(\hat\theta_1,\hat\theta_2)=c,
\]
则
\[
  \Var(\hat\theta_a)
  =
  a^2v_1+(1-a)^2v_2+2a(1-a)c.
\]

\begin{proofbox}[最优权重推导]
对 $a$ 求导：
\[
\begin{aligned}
  \frac{d}{da}\Var(\hat\theta_a)
  &=
  2av_1-2(1-a)v_2+2(1-2a)c.
\end{aligned}
\]
令导数为零并整理：
\[
  a(v_1+v_2-2c)=v_2-c.
\]
因此
\[
  \boxed{
  a^*
  =
  \frac{v_2-c}{v_1+v_2-2c}.
  }
\]
当两个估计量独立时，$c=0$：
\[
  \boxed{
  a^*=\frac{v_2}{v_1+v_2},
  \qquad
  1-a^*=\frac{v_1}{v_1+v_2}.
  }
\]
\end{proofbox}

\begin{intuitionbox}[为什么 $\hat\theta_1$ 的权重公式中出现 $v_2$]
表面上，$\hat\theta_1$ 的权重
\[
  a^*=\frac{v_2}{v_1+v_2}
\]
似乎使用了“另一个估计量”的方差。把它改写为
\[
  a^*
  =
  \frac{1/v_1}{1/v_1+1/v_2}
\]
即可看出：权重与自身方差的倒数成正比。方差越小，精度越高，权重越大。
\end{intuitionbox}

独立时，最小方差为
\[
  \boxed{
  \Var(\hat\theta_{a^*})
  =
  \frac{v_1v_2}{v_1+v_2}.
  }
\]
它小于 $v_1$ 和 $v_2$ 中任一个正数，因此合理合并独立信息可以提高精度。

\begin{examplebox}[最优合并的完整计算]
设两个独立无偏估计量满足
\[
  \Var(T_1)=\frac4n,
  \qquad
  \Var(T_2)=\frac9n.
\]
对组合
\[
  T_a=aT_1+(1-a)T_2,
\]
最优权重为
\[
  a^*
  =
  \frac{9/n}{4/n+9/n}
  =
  \frac9{13}.
\]
因此
\[
  T_{a^*}
  =
  \frac9{13}T_1+\frac4{13}T_2.
\]
$T_1$ 方差更小，所以获得更大权重。最小方差为
\[
  \Var(T_{a^*})
  =
  \frac{(4/n)(9/n)}{4/n+9/n}
  =
  \boxed{\frac{36}{13n}}.
\]
\end{examplebox}

\subsection{样本均值为何无偏}

设 $X_1,\ldots,X_n$ 独立同分布，且
\[
  \E(X_i)=\mu,
  \qquad
  \Var(X_i)=\sigma^2.
\]
样本均值为
\[
  \bar X=\frac1n\sum_{i=1}^nX_i.
\]
利用期望的线性性质：
\[
\begin{aligned}
  \E(\bar X)
  &=
  \E\left(\frac1n\sum_{i=1}^nX_i\right)\\
  &=
  \frac1n\sum_{i=1}^n\E(X_i)\\
  &=
  \frac1n(n\mu)\\
  &=
  \boxed{\mu}.
\end{aligned}
\]
所以 $\bar X$ 是 $\mu$ 的无偏估计量。

独立性用于计算方差：
\[
\begin{aligned}
  \Var(\bar X)
  &=
  \Var\left(\frac1n\sum_{i=1}^nX_i\right)\\
  &=
  \frac1{n^2}\sum_{i=1}^n\Var(X_i)\\
  &=
  \frac1{n^2}(n\sigma^2)\\
  &=
  \boxed{\frac{\sigma^2}{n}}.
\end{aligned}
\]

\begin{conceptbox}[期望计算与方差计算对独立性的需求不同]
期望的线性性质总是成立，因此证明 $\E(\bar X)=\mu$ 不需要独立性；但要把和的方差写成各方差之和，通常需要独立或至少协方差为零。
\end{conceptbox}

\subsection{样本方差为何除以 \texorpdfstring{$n-1$}{n-1}}

先定义使用分母 $n$ 的量：
\[
  S_n^2
  =
  \frac1n\sum_{i=1}^n(X_i-\bar X)^2.
\]
它看起来很自然，但它对总体方差 $\sigma^2$ 有向下偏差。无偏样本方差定义为
\[
  \boxed{
  S^2
  =
  \frac1{n-1}\sum_{i=1}^n(X_i-\bar X)^2.
  }
\]

\begin{intuitionbox}[自由度直觉]
残差 $X_i-\bar X$ 满足
\[
  \sum_{i=1}^n(X_i-\bar X)=0.
\]
因此知道前 $n-1$ 个残差后，最后一个残差已经被上述约束决定。虽然表面上有 $n$ 个残差，真正可以自由变化的只有 $n-1$ 个，所以用于估计方差的自由度为 $n-1$。
\end{intuitionbox}

\begin{proofbox}[关键恒等式]
对每个 $i$，
\[
  X_i-\bar X
  =
  (X_i-\mu)-(\bar X-\mu).
\]
平方求和：
\[
\begin{aligned}
  \sum_{i=1}^n(X_i-\bar X)^2
  &=
  \sum_{i=1}^n
  \left[(X_i-\mu)-(\bar X-\mu)\right]^2\\
  &=
  \sum_{i=1}^n(X_i-\mu)^2
  {}+n(\bar X-\mu)^2\\
  &\quad
  -2(\bar X-\mu)\sum_{i=1}^n(X_i-\mu).
\end{aligned}
\]
注意
\[
  \sum_{i=1}^n(X_i-\mu)
  =
  n(\bar X-\mu).
\]
所以
\[
  \boxed{
  \sum_{i=1}^n(X_i-\bar X)^2
  =
  \sum_{i=1}^n(X_i-\mu)^2
  -
  n(\bar X-\mu)^2.
  }
\]
\end{proofbox}

对恒等式两边取期望：
\[
\begin{aligned}
  \E\left[
  \sum_{i=1}^n(X_i-\bar X)^2
  \right]
  &=
  \sum_{i=1}^n\E[(X_i-\mu)^2]
  -
  n\E[(\bar X-\mu)^2]\\
  &=
  n\sigma^2
  -
  n\Var(\bar X)\\
  &=
  n\sigma^2
  -
  n\left(\frac{\sigma^2}{n}\right)\\
  &=
  (n-1)\sigma^2.
\end{aligned}
\]
因此
\[
  \E(S_n^2)
  =
  \frac{n-1}{n}\sigma^2,
\]
而
\[
  \boxed{
  \E(S^2)
  =
  \frac1{n-1}(n-1)\sigma^2
  =
  \sigma^2.
  }
\]

\begin{intuitionbox}[为什么用样本均值会低估波动]
样本均值 $\bar X$ 是专门根据当前样本选择出来的中心，它使
\[
  \sum(X_i-\bar X)^2
\]
达到最小。因此样本点围绕 $\bar X$ 的平方距离，通常比它们围绕真实但未知的 $\mu$ 的平方距离更小。除以 $n-1$ 正是在平均意义上修正这一低估。
\end{intuitionbox}

\subsection{\texorpdfstring{$S^2$ 无偏不代表 $S$ 无偏}{S2 无偏不代表 S 无偏}}

已知
\[
  \E(S^2)=\sigma^2.
\]
但通常不能对期望直接开平方：
\[
  \E(S)
  =
  \E(\sqrt{S^2})
  \ne
  \sqrt{\E(S^2)}
  =
  \sigma.
\]
由于平方根函数是凹函数，Jensen 不等式给出
\[
  \E(S)\leq\sqrt{\E(S^2)}=\sigma,
\]
在非退化情形下通常严格小于。因此样本标准差 $S$ 通常向下有偏。

\begin{warnbox}[常见错误]
\begin{itemize}
  \item 从 $\E(T)=\theta$ 不能推出 $\E[g(T)]=g(\theta)$；
  \item 从 $S^2$ 对 $\sigma^2$ 无偏，不能推出 $S$ 对 $\sigma$ 无偏；
  \item 从两个估计量都无偏，不能推出它们方差相同；
  \item 比较估计量时，不能只说“这个无偏，所以一定更好”；
  \item 计算线性组合方差时，不能忘记协方差项。
\end{itemize}
\end{warnbox}

\subsection{综合例题}

\begin{examplebox}[例题：比较两个估计量]
设 $X_1,\ldots,X_n$ 独立同分布，均值为 $\mu$、方差为 $\sigma^2$。比较
\[
  T_1=\bar X,
  \qquad
  T_2=\frac{n-1}{n}\bar X.
\]

\textbf{第一步：计算偏差。}
\[
  \Bias(T_1)=0,
\]
\[
  \Bias(T_2)
  =
  \frac{n-1}{n}\mu-\mu
  =
  -\frac{\mu}{n}.
\]

\textbf{第二步：计算方差。}
\[
  \Var(T_1)=\frac{\sigma^2}{n},
\]
\[
  \Var(T_2)
  =
  \left(\frac{n-1}{n}\right)^2\frac{\sigma^2}{n}
  =
  \frac{(n-1)^2\sigma^2}{n^3}.
\]

\textbf{第三步：计算 MSE。}
\[
  \MSE(T_1)=\frac{\sigma^2}{n},
\]
\[
  \MSE(T_2)
  =
  \frac{(n-1)^2\sigma^2}{n^3}
  +
  \frac{\mu^2}{n^2}.
\]

$T_2$ 虽然方差更小，但有偏。因此不能仅凭无偏性或方差判断谁更好；需要比较上述两个 MSE，而比较结果可能依赖 $\mu$ 与 $\sigma^2$。
\end{examplebox}

\subsection{答题模板}

\begin{methodbox}[证明某估计量无偏]
\[
  \E(\hat\theta)
  =
  \text{逐步使用期望线性性质}
  =
  \theta.
\]
最后明确写出：“因此 $\hat\theta$ 是 $\theta$ 的无偏估计量。”
\end{methodbox}

\begin{methodbox}[计算某估计量的 MSE]
\begin{enumerate}
  \item 计算 $\E(\hat\theta)$；
  \item 得到 $\Bias(\hat\theta)=\E(\hat\theta)-\theta$；
  \item 计算 $\Var(\hat\theta)$；
  \item 使用
  \[
    \MSE(\hat\theta)
    =
    \Var(\hat\theta)+\Bias(\hat\theta)^2.
  \]
\end{enumerate}
\end{methodbox}

\begin{methodbox}[合并两个无偏估计量]
\begin{enumerate}
  \item 先验证权重和为 1，从而组合仍无偏；
  \item 写出完整方差，包含协方差项；
  \item 对权重求导并令其为零；
  \item 检查所得权重确实给出最小值；
  \item 解释方差较小的估计量为何获得更大权重。
\end{enumerate}
\end{methodbox}

\begin{summarybox}[第 2 部分的核心]
\[
  \boxed{
  \text{Bias 描述中心位置，Variance 描述波动，MSE 同时衡量二者。}
  }
\]
\[
  \boxed{
  \MSE(\hat\theta)
  =
  \Var(\hat\theta)+\Bias(\hat\theta)^2.
  }
\]
无偏估计量的 MSE 等于方差，但有偏估计量也可能通过降低方差取得更小的 MSE。样本均值是总体均值的无偏估计；样本方差必须除以 $n-1$，才能修正使用样本均值作为中心造成的平均低估。
\end{summarybox}

\section{最后自测}

\begin{enumerate}
  \item 参数与估计量中，哪一个在频率学派中固定，哪一个随机？
  \item 标准差与标准误分别描述什么？
  \item 为什么配对样本不能直接当作两个独立样本？
  \item 为什么总体方差未知时，单样本均值推断使用 $t$ 而不是 $Z$？
  \item 无偏估计量是否一定比有偏估计量具有更小的 MSE？
  \item 在 MSE 分解证明中，交叉项为什么为零？
  \item 为什么独立估计量的最优组合权重与方差的倒数成正比？
  \item 为什么样本方差使用 $n-1$？这与自由度有什么关系？
  \item 为什么 $S^2$ 无偏不能推出 $S$ 无偏？
\end{enumerate}

\begin{goalbox}[自测答案提示]
\begin{enumerate}
  \item 参数固定未知；估计量随样本变化。
  \item 标准差描述个体波动；标准误描述估计量的抽样波动。
  \item 配对数据存在一一对应关系，应分析差值。
  \item 用随机的 $S$ 替换 $\sigma$ 引入额外不确定性，得到 $t$ 分布。
  \item 不一定，应比较 MSE。
  \item 因为 $\E[\hat\theta-\E(\hat\theta)]=0$。
  \item 方差越小代表精度越高，所以应获得更大权重。
  \item 估计样本均值消耗一个自由度，且围绕 $\bar X$ 的平方距离平均偏小。
  \item 一般有 $\E[g(T)]\ne g(\E T)$；平方根还是凹函数。
\end{enumerate}
\end{goalbox}

\end{document}
